I have written the estimated time for each part in parentheses under subsection label. Generally, setting up diy and opam, as well as relearning a bit of C, made up the lion's share of the time.

\section{} \label{}
% details on diy7 etc 
%https://diy.inria.fr/doc/gen.html#diycross%3Aintro 
\subsection{(2 hours)}
We use these commands to generate the tests:
\begin{lstlisting}
    diyone7 -arch AArch64 "PodWR Fre PodWR Fre" -name SB
    diyone7 -arch AArch64 "PodWR Fre PodWR Fre PodWR Fre" -name SB3    
    diyone7 -arch AArch64 "DMB.SYdWR Fre DMB.SYdWR Fre" -name SB+fen 
    diyone7 -arch AArch64 "DMB.SYdWR Fre DMB.SYdWR Fre DMB.SYdWR Fre" -name SB3+fen
\end{lstlisting}

'SB3' is my shorthand for Store Buffering with 3 threads. 
I am on an M4 Max chip (AArch64), and I expect to observe significant relaxations in the unfenced SB and SB3 tests.

\lstinputlisting[caption={System Specs}, label={lst:specs}]{../logs/q1/sys_summary.txt}

\subsection{(30 min)}

Running with litmus7, we notice 7.7 and 9.4 seconds to run SB with and without barrier respectively. 
As expecteed, the test condition is not validated in the fenced case, but is in the unfenced case.

For SB3, we see 10.4 and 12.4 seconds respectively, a slightly higher percentage difference (81\% vs 83\%). 
We see the expected fenced vs unfenced results. 

Another curious difference is that while unfenced SB has only 234 out of 1e8 iterations validating the test condition, SB3 has 1870409 out of 1e8 iterations validating the test condition.
We may see more prevalent compiler optimization and reordering in the latter case with more threads, 
possibly the extra thread results in more optimization opportunities.

Using -s 1k, -r 100k to get 1e8 for each case:

\begin{lstlisting}
    litmus7 SB.litmus -s 1k -r 100k > logs/q1/b_1.txt 
    litmus7 SB+fen.litmus -s 1k -r 100k > logs/q1/b_1fen.txt     
    litmus7 SB3.litmus -s 1k -r 100k > logs/q1/b_2.txt 
    litmus7 SB3+fen.litmus -s 1k -r 100k > logs/q1/b_2fen.txt 
\end{lstlisting}


\subsection{(2 min)}

Comparing the axiomatic results, we naturally notice no difference in timing. 
However we notice the expected difference in the results, with our test condition missing from possible outputs in the fenced cases. 

\begin{lstlisting}
    herd7 SB.litmus > logs/q1/c_1.txt 
    herd7 SB3.litmus > logs/q1/c_1fen.txt  
    herd7 SB+fen.litmus > logs/q1/c_2.txt
    herd7 SB3+fen.litmus > logs/q1/c_2fen.txt
\end{lstlisting}
%https://www.cl.cam.ac.uk/~sf502/regressions/rmem/help.html

\subsection{(30 min)}
We use interactive -false and eager -true flags. 
In both cases, the fenced version notes 'allowed not found'  agreeing with prior analaysis. 
\begin{lstlisting}    
    rmem -model flat -interactive false -eager true SB.litmus > logs/q1/d_1.txt
    rmem -model flat -interactive false -eager true SB+fen.litmus > logs/q1/d_1fen.txt
    rmem -model flat -interactive false -eager true SB3.litmus > logs/q1/d_2.txt      
    rmem -model flat -interactive false -eager true SB3+fen.litmus > logs/q1/d_2fen.txt
\end{lstlisting}

\subsection{(20 min)}
Again, we reinforced the expected results from the above. 
\begin{lstlisting}
    mcompare7 logs/q1/b_1.txt logs/q1/c_1.txt > logs/q1/e_1.txt
    mcompare7 logs/q1/b_1fen.txt logs/q1/c_1fen.txt > logs/q1/e_1fen.txt
    mcompare7 logs/q1/b_2.txt logs/q1/c_2.txt > logs/q1/e_2.txt
    mcompare7 logs/q1/b_2fen.txt logs/q1/c_2fen.txt > logs/q1/e_2fen.txt
    
\end{lstlisting}


%LINK: 
% https://www.cl.cam.ac.uk/~sf502/regressions/rmem/
